% 问题四正式正文：稀疏短期模型评测数据下的前沿能力预测。
% 图件来自 paper/figures/q4/，由 scripts/plot_q4_figures.py 生成。

\section{问题四：稀疏短期模型评测数据下的前沿能力预测}
\label{sec:q4}

\subsection{问题分析与研究目标}

问题四研究的是一个稀疏、短期且横截面高度不均衡的预测问题：在可比评测月份较少、每月模型提交数量较多的条件下，估计模型能力前沿在未来十二个月的可能范围。与单一模型的时间序列预测不同，本文的目标变量是每个月模型横截面综合分数的高分位数，因此需要同时处理评分口径、模型构成和时间外推的不确定性。

本文的核心主张限定为：在当前附件 C 的可比窗口内，详细 C8 六任务等权分数可以支持一个低复杂度、可复核的上尾预测框架；该框架能够比较平台型和趋势型情景，但不能据此识别纯技术进步的因果效应，也不能给出统计上精确的长期点预测。问题四的数据建模不依赖问题一至问题三的中间结果；问题三的规模弹性只作为算力情景层的外生输入。

\subsection{数据边界、投毒隔离与分析流程}

研究数据来自附件 C 的 \texttt{C\_efficiency\_evolution} 目录。虽然扩展时间序列覆盖 2019--2025 年，但 2019--2023 年仅有少量来自不同来源的历史记录，缺少与近年榜单一致的任务定义和缺失处理规则。因此主分析限定为同一 Open LLM Leaderboard 来源下的 2024 年 6 月至 2025 年 3 月，共十个月。题面中预填的系数、相关性、预测结果和提示性答案不被当作观测值或模型参数；仅保留可追溯的榜单字段和详细 C8 任务结果。

数据处理先完成字段审计、模型名称连接、日期解析、重复检查和缺失记录计数，再进行 C8 重算。榜单与 C8 详细结果合并后，1,892 行中保留 1,889 行，排除 3 行缺少有效日期或 C8 分数的记录；官方 Average 表的 4,576 行中保留 4,564 行，排除 12 行。所有排除数量均写入图件源数据目录，未为绘图便利进行抽样。

图~\ref{fig:q4-workflow} 给出从数据审计到预测情景的完整流程。问题二和问题三的结果在流程末端以外生约束进入情景投影，不回写历史观测，也不替代问题四自身的时间外验证。

\begin{figure}[htbp]
  \centering
  \includegraphics[width=0.98\textwidth]{FigQ4_03_workflow.pdf}
  \caption{问题四分析流程。流程依次经过投毒隔离、C8 等权聚合、月度前沿计算、70/30 时间外验证、分位数模型和十二个月情景投影。问题二、问题三的输入只进入情景层。}
  \label{fig:q4-workflow}
\end{figure}

\subsection{指标定义与预测模型}

主目标由详细 C8 JSON 中的 IFEval、BBH、MATH、GPQA、MUSR 和 MMLU-PRO 六个一级任务构造。对模型 $i$，六项有效任务分数等权平均并转换到 0--100 分尺度：

\begin{equation}
 S_i=100\times\frac{1}{6}\sum_{k=1}^{6}s_{ik}.
 \label{eq:q4-c8}
\end{equation}

按月份 $t$ 聚合模型横截面，详细 C8 前沿定义为

\begin{equation}
 F_t=Q_{0.95}\left(S_i\mid \mathrm{month}=t\right).
 \label{eq:q4-frontier}
\end{equation}

榜单官方 \texttt{Average} 保留为敏感性口径，不能与详细 C8 混合平均，也不能用一种口径的预测器评价另一种口径的目标。

滚动基线从至少三个月历史开始，比较最近值、近三月均值、扩展线性趋势和最近四点线性趋势。一步预测报告 MAE、RMSE、Bias 和普通 $R^2$；$h=3$ 只在测试起点足够时报告描述性误差。对于面板数据，本文拟合低复杂度的条件分位数模型：

\begin{equation}
 Q_{\tau}(S_i\mid X_i)=\alpha_{\tau}+\beta_{\tau}\log(1+N_i)+\gamma_{\tau}t_i+\delta_{\tau,\mathrm{type}(i)},
 \qquad \tau\in\{0.90,0.95\},
 \label{eq:q4-quantile}
\end{equation}

其中 $N_i$ 为参数规模，$t_i$ 为相对起始月份，\texttt{type} 为模型类型。模型选择依据主要为分位数 pinball loss、经验覆盖率和时间外稳定性，普通 $R^2$ 仅作辅助诊断。

\subsection{实验协议与不确定性}

主协议按月份进行 70/30 时间切分：2024 年 6 月至 12 月为训练段，2025 年 1 月至 3 月为测试段。全类型实验包含 1,167 个训练模型和 684 个测试模型；pretrained-only 结果只作类型敏感性分析。月度前沿的阴影区间来自已保存的月度横截面 bootstrap 95\% 区间，而非来自新增的合成观测。

未来十二个月情景采用问题三在 $L_{\mathrm{ctx}}=4096$ 条件下得到的规模弹性作为外生输入：

\begin{equation}
 N^*(C)\propto C^{0.452},\qquad D^*(C)\propto C^{0.548}.
 \label{eq:q4-compute}
\end{equation}

在末月模型类型构成下考察 1.0、1.5、2.3、4.0 和 5.0 倍年算力倍率。平台型情景固定时间项，趋势型情景将时间项推进十二个月；类型比例的 $\pm10$ 个百分点扰动用于敏感性分析。profile range 只表示模型类型重加权范围，不是时间预测置信区间。

\subsection{结果：数据趋势、短期验证与任务异质性}

图~\ref{fig:q4-frontier} 汇总了当前主协议的四类证据。月度详细 C8 q95 从 2024 年 6 月的约 45.8 分升至 2025 年 3 月的约 54.2 分；官方 Average q95 处于另一尺度，说明目标定义本身是预测敏感性来源。任务级变化中，MATH 的首末月 q95 增加 27.30 分，IFEval、BBH、MMLU-PRO 和 GPQA 分别增加 5.54、3.35、4.86 和 2.55 分，MUSR 基本不变。排除 MATH 后综合前沿仍增加约 4.2 分，因此综合趋势被 MATH 放大，但并非完全由单一任务造成。

滚动验证显示，全类型 C8 的 $h=1$ 最佳方法为最近四点线性趋势，MAE=0.852、RMSE=0.860、普通 $R^2=0.700$，但测试点仅有 3 个；$h=3$ 只有 1 个测试起点。pretrained-only 的 $h=1$ MAE 为 9.924--13.503，表现明显更不稳定。图中模型类型构成也随月份变化，说明时间项可能同时吸收技术进步、提交者构成和评测版本变化。

\begin{figure}[htbp]
  \centering
  \includegraphics[width=0.98\textwidth]{FigQ4_01_frontier_validation.pdf}
  \caption{前沿数据、任务异质性与短期验证。（a）详细 C8 与官方 Average 的月度 q95，阴影为月度横截面 bootstrap 95\% 区间，浅色区域为测试月份；（b）六个任务首末月 q95 变化；（c）前沿样本的模型类型构成；（d）$h=1$ 滚动验证的 MAE 与普通测试 $R^2$。$h=1$ 测试点数为 3，$h=3$ 仅有 1 个起点。}
  \label{fig:q4-frontier}
\end{figure}

70/30 截面 q95 模型的全类型 pinball loss 为 0.559，普通测试 $R^2=-0.425$；pretrained-only 的 pinball loss 为 0.407，普通测试 $R^2=-1.139$。由于模型目标是条件上尾，pinball loss 和覆盖率比普通 $R^2$ 更适合作为主评价。早期年度切分的覆盖率约为 95\%，但只作为兼容性检查，不与当前主协议重复计数。

\subsection{结果：十二个月算力—时间情景}

图~\ref{fig:q4-scenarios} 和表~\ref{tab:q4-scenarios} 给出基准类型构成下的十二个月情景。1.0、2.3 和 5.0 倍算力下，平台型前沿分别为 58.18、61.28 和 64.20 分，趋势型前沿分别为 69.68、72.78 和 75.71 分。平台型与趋势型之间的差异大于单独提高算力倍率带来的差异，说明时间项外推是当前结果的主要结构性敏感性。

\begin{figure}[htbp]
  \centering
  \includegraphics[width=0.98\textwidth]{FigQ4_02_12m_scenarios.pdf}
  \caption{未来十二个月的算力—时间情景。（a）不同年算力倍率下平台型和趋势型 C8 q95；（b）固定 2.3 倍算力时，对 pretrained 或 chat 类型比例进行 $\pm10$ 个百分点扰动后的加权平均分。profile range 是类型构成重加权范围，不是时间预测置信区间。}
  \label{fig:q4-scenarios}
\end{figure}

\begin{table}[htbp]
\centering
\caption{详细 C8 口径下的十二个月算力与时间情景}
\label{tab:q4-scenarios}
\renewcommand{\arraystretch}{1.15}
\begin{tabular}{ccc}
\toprule
年算力倍率 & 平台型 & 趋势型\\
\midrule
1.0$\times$ & 58.18 & 69.68\\
1.5$\times$ & 59.68 & 71.19\\
2.3$\times$ & 61.28 & 72.78\\
4.0$\times$ & 63.36 & 74.86\\
5.0$\times$ & 64.20 & 75.71\\
\bottomrule
\end{tabular}
\end{table}

\subsection{讨论：证伪结果与适用边界}

上述结果不支持把时间系数直接解释为纯技术进步。时间项可能同时包含模型类型比例、评测版本、提交者组成和数据源漂移；而主协议只有十个月可比窗口和三个月测试段，无法将这些因素分离。与此同时，High 组 Loss--Benchmark 桥接只有 7 个可比样本，斜率区间跨过零，Medium 组训练数据量缺失较多，因此该桥接不能替代详细 C8 直接前沿模型。

问题四的证据边界有三点。第一，$h=1$ 的测试点数为 3、$h=3$ 的测试点数为 1，不能支撑复杂模型选择或显著性检验。第二，基础模型族、训练数据量、训练算力和评测版本字段不完整，发布方前缀只能作为组织聚类代理。第三，算力情景使用问题三的规模弹性，而不是同一模型族上的训练轨迹，因此属于条件情景而非因果估计。基于这些边界，正式结果应并列报告平台型和趋势型范围，不应报告单一精确的长期点预测。

\subsection{结论}

本文为稀疏短期模型评测数据建立了一个可复现的详细 C8 上尾预测流程：先隔离题面投毒和不可追溯结果，再计算六任务等权月度前沿，使用滚动基线和 70/30 时间外面板分位数模型进行验证，最后在问题三规模弹性约束下输出十二个月情景。实验支持在当前数据范围内比较不同时间项和算力倍率的相对影响，但不支持因果解释或高置信度长期点预测。该框架的可复用价值在于把评分口径、模型构成、短样本误差和情景外生条件放在同一审计链条中报告。
